\section{Задание}
Реализовать полносвязную нейронную сеть для распознавания рукописных цифр MNIST. Самостоятельно прочитать данные IDX, реализовать прямой и обратный проход, вычисление функции потерь и обновление параметров методом Adam. Сравнить архитектуры по точности на валидационной выборке, оценить выбранную модель на тесте и проверить её на внешних рисунках.

\section{Теоретические сведения}
Слой \texttt{Flatten} преобразует изображение \(28\times28\) в вектор из 784 элементов. Для батча \(X\) слой \texttt{Linear} вычисляет \(Y=XW+b\), где \(W\)~--- матрица весов, а \(b\)~--- вектор смещений. Для градиента \(G=\partial L/\partial Y\) обратный проход имеет вид
\[
    \frac{\partial L}{\partial X}=GW^\top,\qquad
    \frac{\partial L}{\partial W}=X^\top G,\qquad
    \frac{\partial L}{\partial b}=\sum_i G_{i,:}.
\]
После скрытого линейного слоя применяется \(\operatorname{ReLU}(x)=\max(0,x)\). При обратном проходе она пропускает градиент только для положительных входов. Скрытые веса инициализируются с дисперсией \(2/n_{\mathrm{in}}\) по Хе, выходные~--- с дисперсией \(1/n_{\mathrm{in}}\). Смещения равны нулю в начале обучения.

Для логитов \(z_{ik}\), где \(i\)~--- номер примера, \(k\)~--- класс, положим \(\tilde z_{ik}=z_{ik}-\max_j z_{ij}\). Тогда
\[
  p_{ik}=\frac{\exp(\tilde z_{ik})}{\sum_j\exp(\tilde z_{ij})},\qquad
  L=\frac1N\sum_i\left(\log\sum_j\exp(\tilde z_{ij})-\tilde z_{i,y_i}\right),
  \qquad
  \frac{\partial L}{\partial z_{ik}}=\frac{p_{ik}-\mathbf 1[k=y_i]}N.
\]
Вычитание максимума защищает экспоненту от переполнения. \texttt{SoftmaxWithLoss} хранит вероятности и one-hot-метки для вычисления градиента.

В режиме обучения Dropout сохраняет каждую активацию с вероятностью \(1-p\) и делит сохранённые значения на \(1-p\), сохраняя их математическое ожидание. В режиме оценки слой пропускает вход без преобразований.

Adam на шаге \(t\) обновляет скользящие средние градиента \(g_t\) и его квадрата. Для каждого обучаемого параметра \(\theta\) вычисляются
\begin{align*}
 m_t&=\beta_1 m_{t-1}+(1-\beta_1)g_t,&
 v_t&=\beta_2 v_{t-1}+(1-\beta_2)g_t^2,\\
 \hat m_t&=m_t/(1-\beta_1^t),&
 \hat v_t&=v_t/(1-\beta_2^t),&
 \theta_t&=\theta_{t-1}-\alpha\frac{\hat m_t}{\sqrt{\hat v_t}+\varepsilon}.
\end{align*}

\clearpage
\section{Подготовка данных}
Загрузчик читает заголовки IDX в порядке big-endian, проверяет магические числа 2051 и 2049, размерности и длину массива байтов. Изображения образуют массивы \((60\,000,28,28)\) и \((10\,000,28,28)\), метки~--- векторы соответствующих длин. Пиксели преобразуются в \texttt{float32} и делятся на 255, метки кодируются в one-hot формы \((N,10)\).

При seed 42 индексы исходной обучающей выборки перемешиваются. Первые 6\,000 объектов используются для валидации, остальные 54\,000~--- для обучения. Разбиение одинаково для всех моделей. Тестовый набор используется после выбора конфигурации и эпохи.

\begin{table}[!ht]
\centering\small
\caption{Распределение изображений по классам}
\begin{tabular}{@{}lrrrrrrrrrrr@{}}
\toprule
Выборка & 0 & 1 & 2 & 3 & 4 & 5 & 6 & 7 & 8 & 9 & Всего\\
\midrule
Исходная & 5923 & 6742 & 5958 & 6131 & 5842 & 5421 & 5918 & 6265 & 5851 & 5949 & 60000\\
Train & 5319 & 6067 & 5361 & 5487 & 5255 & 4913 & 5329 & 5657 & 5243 & 5369 & 54000\\
Validation & 604 & 675 & 597 & 644 & 587 & 508 & 589 & 608 & 608 & 580 & 6000\\
Test & 980 & 1135 & 1032 & 1010 & 982 & 892 & 958 & 1028 & 974 & 1009 & 10000\\
\bottomrule
\end{tabular}
\end{table}
\experimentfigure{0.59}{mnist_samples.png}{Первые 25 изображений исходного обучающего набора}

\clearpage
\section{Реализация и обучение}
\subsection{Слои сети}
Linear хранит вход, ReLU~--- маску положительных значений, Flatten~--- исходную форму батча, Dropout~--- масштабированную маску. Обратный проход использует состояние последнего обучающего прохода.
{\codefont
\begin{verbatim}
class Linear:
    def __init__(self, n_in, n_out, rng, *, relu=True, dtype=np.float32):
        # He для скрытого слоя с ReLU, масштаб 1/sqrt(n_in) для логитов.
        scale = np.sqrt((2 if relu else 1) / n_in)
        self.W = rng.normal(0, scale, (n_in, n_out)).astype(dtype)
        self.b = np.zeros(n_out, dtype=dtype)
        self.dW = np.zeros_like(self.W)
        self.db = np.zeros_like(self.b)

    def forward(self, x, training=True):
        if training:
            self.x = x
        return x @ self.W + self.b

    def backward(self, dy):
        self.dW[:] = self.x.T @ dy
        self.db[:] = dy.sum(axis=0)
        return dy @ self.W.T

    def parameters(self):
        return ((self.W, self.dW), (self.b, self.db))


class ReLU:
    def forward(self, x, training=True):
        if training:
            self.positive = x > 0
        return np.maximum(x, 0)

    def backward(self, dy):
        return dy * self.positive


class Flatten:
    def forward(self, x, training=True):
        if training:
            self.shape = x.shape
        return x.reshape(len(x), -1)

    def backward(self, dy):
        return dy.reshape(self.shape)


class Dropout:
    def __init__(self, probability, rng):
        if not 0 <= probability < 1:
            raise ValueError("Dropout probability must be in [0, 1).")
        self.keep = 1 - probability
        self.rng = rng

    def forward(self, x, training=True):
        if not training:
            return x
        self.scale = (self.rng.random(x.shape, dtype=np.float32) < self.keep).astype(x.dtype)
        self.scale /= self.keep
        return x * self.scale

    def backward(self, dy):
        return dy * self.scale
\end{verbatim}
}

\clearpage
\subsection{Функция потерь и оптимизатор}
Кросс-энтропия вычисляется по сдвинутым логитам. Adam обновляет веса и смещения на месте после вычисления градиентов всех слоёв.
{\codefont
\begin{verbatim}
def _softmax_terms(logits):
    # Вычитание максимума защищает exp от переполнения.
    shifted = logits - logits.max(axis=1, keepdims=True)
    exp = np.exp(shifted)
    total = exp.sum(axis=1, keepdims=True)
    return shifted, exp / total, total


def softmax(logits):
    return _softmax_terms(logits)[1]


class SoftmaxWithLoss:
    def forward(self, logits, one_hot):
        shifted, self.probabilities, total = _softmax_terms(logits)
        self.targets = one_hot
        # log-sum-exp сохраняет конечный loss даже при exp(z) = 0.
        return float((np.log(total[:, 0]) - (one_hot * shifted).sum(axis=1)).mean())

    def backward(self):
        return (self.probabilities - self.targets) / len(self.targets)


class Adam:
    def __init__(self, parameters, lr=0.001, beta1=0.9, beta2=0.999, eps=1e-8):
        self.parameters = list(parameters)
        self.lr, self.beta1, self.beta2, self.eps = lr, beta1, beta2, eps
        self.m = [np.zeros_like(value) for value, _ in self.parameters]
        self.v = [np.zeros_like(value) for value, _ in self.parameters]
        self.t = 0

    def step(self):
        self.t += 1
        first_correction = 1 - self.beta1**self.t
        second_correction = 1 - self.beta2**self.t
        for (value, gradient), m, v in zip(self.parameters, self.m, self.v):
            m *= self.beta1
            m += (1 - self.beta1) * gradient
            v *= self.beta2
            v += (1 - self.beta2) * np.square(gradient)
            m_hat = m / first_correction
            v_hat = v / second_correction
            value -= self.lr * m_hat / (np.sqrt(v_hat) + self.eps)

\end{verbatim}
}

\subsection{Проверка градиентов}
Для Linear \(3\to4\) с двумя входными векторами аналитические значения \(dX\), \(dW\) и \(db\) сравнены с центральными разностями \((L(\theta+h)-L(\theta-h))/(2h)\). При \(h=10^{-5}\) и вычислениях в \texttt{float64} максимальная абсолютная ошибка составила \(2{,}79\cdot10^{-11}\).

\clearpage
\subsection{Цикл обучения}
Для каждой архитектуры выполнено 12 эпох с мини-батчем 256. Веса и маски Dropout используют генератор с seed 42. Порядок примеров перемешивается в каждой эпохе с использованием генератора с seed 43. Параметры Adam: \(\alpha=0{,}001\), \(\beta_1=0{,}9\), \(\beta_2=0{,}999\), \(\varepsilon=10^{-8}\). Модель выполняет прямой проход по слоям в заданном порядке, а обратный~--- в обратном. Метрики считаются в конце эпохи на полных train и validation с отключённым Dropout. При равной валидационной точности выбирается более ранняя эпоха.
{\codefont
\begin{verbatim}
def train(config, train_x, train_y, val_x, val_y, epochs):
    name, widths, dropout = config
    model = MLP(widths, dropout, SEED)
    best_model = MLP(widths, dropout, SEED)
    optimizer = Adam(model.parameters(), lr=LEARNING_RATE)
    objective = SoftmaxWithLoss()
    shuffle = np.random.default_rng(SEED + 1)
    history = []
    best_epoch, best_accuracy = 0, -1.0
    for epoch in range(1, epochs + 1):
        order = shuffle.permutation(len(train_x))
        for start in range(0, len(order), BATCH_SIZE):
            batch = order[start:start + BATCH_SIZE]
            logits = model.forward(train_x[batch], training=True)
            objective.forward(logits, train_y[batch])
            model.backward(objective.backward())
            optimizer.step()

        train_loss, train_acc, _ = evaluate(model, train_x, train_y)
        val_loss, val_acc, _ = evaluate(model, val_x, val_y)
        history.append((name, epoch, train_loss, train_acc, val_loss, val_acc))
        print(f"{name} {epoch:2d}  loss {train_loss:.5f}/{val_loss:.5f}"
              f"  accuracy {train_acc:.4%}/{val_acc:.4%}", flush=True)

        # Strict comparison keeps the earliest checkpoint when accuracy ties.
        if val_acc > best_accuracy:
            best_epoch, best_accuracy = epoch, val_acc
            for (destination, _), (source, _) in zip(best_model.parameters(), model.parameters()):
                destination[...] = source
    return best_model, best_epoch, best_accuracy, history
\end{verbatim}
}

\section{Сравнение архитектур}
Во всех конфигурациях скрытые Linear сопровождаются ReLU. Переход \(A\to B\) меняет ширину первого слоя, \(B\to C\) добавляет второй слой, \(C\to D\) включает Dropout после каждой ReLU.

\begin{table}[!ht]
\centering\small
\caption{Конфигурации и максимальная точность на валидации}
\begin{tabular}{@{}llrrrr@{}}
\toprule
Модель & Скрытые слои & Dropout & Параметры & Эпоха & Точность, \%\\
\midrule
A & 128 & 0 & 101\,770 & 12 & 97,38\\
B & 256 & 0 & 203\,530 & 12 & 97,90\\
C & \(256\to128\) & 0 & 235\,146 & 12 & 98,03\\
D & \(256\to128\) & 0,2 & 235\,146 & 12 & 98,12\\
\bottomrule
\end{tabular}
\end{table}

\clearpage
\experimentfigure{1}{learning_A.png}{Потери и точность модели A по эпохам}
\experimentfigure{1}{learning_B.png}{Потери и точность модели B по эпохам}
При увеличении ширины со 128 до 256 нейронов лучшая точность на валидации выросла на 0,52 п.\,п. К концу обучения модель B достигла меньшей кросс-энтропии и на train, и на validation.

\clearpage
\experimentfigure{1}{learning_C.png}{Потери и точность модели C по эпохам}
\experimentfigure{1}{learning_D.png}{Потери и точность модели D по эпохам}
Второй скрытый слой дал ещё 0,13 п.\,п. У C с 9-й по 12-ю эпоху потери train продолжили снижаться, а потери val выросли с 0,07528 до 0,08240. Это совместимо с ростом уверенности в части ошибочных ответов.

При Dropout 0,2 разрыв train--val в конце обучения уменьшился с 1,85 до 1,57 п.\,п. Валидационная кросс-энтропия уменьшилась до 0,07744. Выигрыш D относительно C составил пять верных ответов на 6\,000 изображений. Для тестовой оценки сохранены веса модели D с эпохи 12, на которой получена максимальная валидационная точность 98,12\,\%.

\clearpage
\section{Тестовая оценка}
Выбранная модель распознала 9\,818 из 10\,000 цифр. Точность составила \(98{,}18\,\%\), средняя кросс-энтропия~--- \(0{,}06197\). Строки матрицы ошибок соответствуют истинным меткам, столбцы~--- предсказанным.
\experimentfigure{0.68}{confusion_matrix.png}{Матрица ошибок на тестовой выборке}
\experimentfigure{1}{test_errors.png}{Примеры ошибок из трёх наиболее частых направлений путаницы}

Чаще всего модель принимала 9 за 4 (14 случаев), 7 за 2 (11) и 2 за 8 (8). У показанных 9 разомкнутый верх и нисходящий штрих напоминают открытую 4. Загиб верхнего штриха и поперечная черта у 7 могут восприниматься как верх и основание 2. У 2 с индексом 3330 нижний росчерк образует овал, а у 1790 толстые сближенные штрихи сужают просветы. Эти особенности могли сделать обе цифры похожими на 8.

\clearpage
\section{Распознавание внешних изображений}
Для проверки использованы три оригинальных цифровых рисунка, построенные по заданным вручную траекториям штрихов. У цифры 2 тонкий тёмный штрих и смещение влево, у 5~--- светлый широкий штрих на тёмном фоне, у 7~--- поперечная черта и смещение вниз и вправо.

Изображение переводится в оттенки серого. По медиане пикселей рамки оценивается фон, при необходимости изображение инвертируется. После вычитания фона область штриха вписывается в \(20\times20\) с сохранением пропорций, переносится на холст \(28\times28\) и сдвигается по центру масс к \((13{,}5;\,13{,}5)\). Интенсивность делится на 255. Softmax возвращает вероятности всех десяти классов.
\experimentfigure{0.64}{external_drawings.png}{Исходные рисунки, прямое масштабирование и полная обработка}

Оба варианта обработки дали три верных ответа. Полная обработка повысила оценки Softmax для всех трёх рисунков. Для смещённой цифры 7 оценка выросла с \(37{,}63\,\%\) до \(99{,}75\,\%\). Все преобразования проверялись вместе. На других рисунках низкий контраст, шум на границах или обрезанные штрихи могут помешать выделению цифры.

\section{Заключение}
В этом эксперименте расширение первого слоя дало больший прирост валидационной точности, чем добавление второго. Dropout 0,2 уменьшил разрыв между train и val, а модель с двумя скрытыми слоями и Dropout показала лучшую валидационную точность. На независимом тесте она распознала \(98{,}18\,\%\) цифр. Сохранённые веса позволяют проверить новый рисунок и получить вероятности классов в ноутбуке без повторного обучения.

